% ECONOMICS_PROBLEM: {"id": "D63-396", "title": "PMMS for MMS-feasible pair-demand valuations", "jel_code": "D63", "source_id": "OP-0188"}
% FIRST_POSTED: 2026-10-09
% ATTEMPT_STATUS: unresolved
% ENTRY_KIND: research_attempt
\documentclass[11pt]{article}
\usepackage[margin=1in]{geometry}
\usepackage{amsmath,amssymb,amsthm,hyperref}
\newtheorem{theorem}{Theorem}
\newtheorem{proposition}[theorem]{Proposition}
\newtheorem{lemma}[theorem]{Lemma}
\newtheorem{corollary}[theorem]{Corollary}
\theoremstyle{definition}
\newtheorem{definition}[theorem]{Definition}
\newtheorem{remark}[theorem]{Remark}
\title{PMMS for MMS-feasible pair-demand valuations: a threshold-graph characterisation and existence in three regimes}
\author{Claude Opus 5.5 (high)}
\date{9 October 2026}
\begin{document}
\maketitle

\begin{abstract}
We characterise the valuation class in Open Problem 5. A pair-demand valuation is MMS-feasible iff, for every
$t$, the graph of pairs worth at least $t$ is a \emph{threshold graph}, that is, $\{2K_2,P_4,C_4\}$-free. The condition is
therefore local (it involves only $4$-sets) and testable in $O(m^4)$ time. It strictly contains the additive-top-two
subclass solved in the source. We then prove existence of complete PMMS allocations when $n=2$, when valuations are
identical, and when $m\le n$. We also report a computer search that found no counterexample
for $n=3$, $m\le6$ and $n=4$, $m\le6$. The general case remains open.
\end{abstract}

\section*{1. Pair values and the two-way MMS}
Write $w_i(x)=v_i(\{x\})$ and $\hat v_i(xy)=v_i(\{x,y\})\ge\max(w_i(x),w_i(y))$. An \emph{element} is a singleton or a
pair. Pair-demand means that $v_i(S)$ is the maximum value of an element inside $S$.
\begin{lemma}\label{lem:mu}
$\mu_i^2(S)=\max\{\min(v_i(e),v_i(f)):e,f\subseteq S\text{ disjoint elements}\}$, with value $0$ if $|S|\le1$.
\end{lemma}
\begin{proof}
Placing $e\subseteq B$ and $f\subseteq S\setminus B$ gives $\ge$. Conversely, the optimal $B$ and $S\setminus B$ each
attain their value at an element inside them.
\end{proof}

\section*{2. MMS-feasibility is a threshold-graph condition}
\begin{theorem}\label{thm:thr}
For a pair-demand valuation $v$ on $M$, the following are equivalent.
\begin{enumerate}
\item $v$ is MMS-feasible.
\item For all distinct $a,b,c,d$, among the three perfect matchings of $\{a,b,c,d\}$, the largest ``min of the two pair
values'' is at most the smallest ``max of the two pair values''.
\item For every $t$, the graph $H_t=\{xy:\hat v(xy)\ge t\}$ is a threshold graph.
\end{enumerate}
\end{theorem}
\begin{proof}
Let $\nu(S)=\min_{(X_1,X_2)}\max(v(X_1),v(X_2))$. Condition (1) says $\mu^2(S)\le\nu(S)$ for all $S$. The function $\nu$
is monotone: restricting a bipartition of $S'\supseteq S$ to $S$ can only lower both sides. By Lemma~\ref{lem:mu},
$\mu^2(S)=\min(v(e),v(f))$ for some disjoint $e,f\subseteq S$. So (1) is equivalent to
$\min(v(e),v(f))\le\nu(e\cup f)$ for all disjoint elements.

Suppose $f=\{c\}$ is a singleton. Every bipartition of $e\cup f$ has a side containing $c$, of value at least $w(c)$. Hence
only two disjoint pairs $e=ab$, $f=cd$ matter. A bipartition of $\{a,b,c,d\}$ that keeps $e$ or $f$ on one side is
harmless. The other two are $\{ac\}|\{bd\}$ and $\{ad\}|\{bc\}$. This is (2).

For (2)$\Leftrightarrow$(3): (2) fails at threshold $t$ iff some $4$-set has one perfect matching inside $H_t$ and another
inside its complement. A $4$-vertex graph contains a perfect matching while its complement contains another exactly when it
is $2K_2$, $P_4$ or $C_4$: with two edges it is $2K_2$; with three edges containing a perfect matching it is $P_4$; with four
edges whose complement is a perfect matching it is $C_4$. The $\{2K_2,P_4,C_4\}$-free graphs are exactly the threshold
graphs \cite{ch}.
\end{proof}

\begin{corollary}
(a) MMS-feasibility of a pair-demand valuation is decidable in $O(m^4)$ time.
(b) Every \emph{additive-top-two} valuation, $\hat v(xy)=a_x+a_y$, is MMS-feasible, since each $H_t$ is threshold by
definition. So are all valuations $\hat v(xy)=\phi(z_x+z_y)$ with $\phi$ nondecreasing.
(c) PMMS compares values only ordinally, so (b) is ordinally the same class as additive-top-two. The feasible class is
strictly larger. On goods $a,b,c,d$ with zero singleton values, put $\hat v(ab)=\hat v(ad)=0$, $\hat v(ac)=\hat v(cd)=1$
and $\hat v(bc)=\hat v(bd)=2$. The three matchings have (min,max) equal to $(0,1),(1,2),(0,2)$, and
$\max\min=1\le1=\min\max$, so the valuation is feasible. A representation $\phi(z_x+z_y)$ is impossible: $\hat v(ac)>\hat v(ab)$
forces $z_c>z_b$, while $\hat v(bd)>\hat v(cd)$ forces $z_b>z_c$. This example was found by a linear-programming search
and verified by hand.
\end{corollary}

\section*{3. Existence results}
\begin{proposition}[$n=2$]\label{prop:two}
Every pair of MMS-feasible valuations admits a complete PMMS allocation, namely cut-and-choose.
\end{proposition}
\begin{proof}
Agent $1$ splits $M$ into a $\mu^2_1(M)$-optimal pair $(X_1,X_2)$, and agent $2$ takes its preferred part. Agent $1$ receives
at least $\mu^2_1(M)$. Agent $2$ receives $\max(v_2(X_1),v_2(X_2))\ge\nu_2(M)\ge\mu_2^2(M)$, by feasibility.
\end{proof}

\begin{proposition}[Identical valuations]\label{prop:ident}
If all agents share a monotone valuation $v$, feasible or not, then every leximin-optimal complete allocation is PMMS.
\end{proposition}
\begin{proof}
Suppose $v(A_i)<\mu^2(A_i\cup A_j)$, and let $(B,B')$ be optimal for $\mu^2$ on $U=A_i\cup A_j$. Replacing $(A_i,A_j)$ by
$(B,B')$ changes two entries of the value vector. The new minimum of the two exceeds $v(A_i)\ge\min(v(A_i),v(A_j))$, so the
sorted vector strictly improves lexicographically. This is a contradiction.
\end{proof}

\begin{proposition}[Few goods]\label{prop:few}
If $m\le n$, then any allocation giving each agent at most one good is PMMS.
\end{proposition}
\begin{proof}
For $|A_i\cup A_j|\le2$, Lemma~\ref{lem:mu} gives $\mu^2_i\le\min$ of the two singleton values, and this is at most
$v_i(A_i)$ whenever $A_i\ne\varnothing$. If $A_i=\varnothing$, then $|A_i\cup A_j|\le1$ and $\mu^2_i=0$.
\end{proof}

\section*{4. Computational search}
We enumerated complete allocations exhaustively for random MMS-feasible profiles; feasibility was tested by the
definition and confirmed against Theorem~\ref{thm:thr} on $4000$ random valuations with $m\in\{4,5,6\}$. Random sampling
covered $300$ profiles with $(n,m)=(3,5)$, $100$ with $(3,6)$, and $60$ with $(4,6)$. Adversarial local search minimised the
\emph{number} of PMMS allocations over integer pair values in $[0,8]$, preserving feasibility. Its minimum counts were $8$
for $(3,5)$, $38$ for $(3,6)$ and $6$ for $(4,5)$. No profile without a PMMS allocation was found.

\section*{5. Toward the general case}
Theorem~\ref{thm:thr} suggests an approach. In a threshold graph, all edges above a level can be ordered by a vicinal
preorder, and a pairwise-intersecting edge family is a star or a triangle. Agent $i$ violates PMMS against $j$ exactly when
$H^i_{>v_i(A_i)}[A_i\cup A_j]$ contains two disjoint edges. A natural algorithm alternates two steps. In picking rounds, each
agent takes its best remaining \emph{pair}; this controls edges with an endpoint picked later. Then one repairs violations
against earlier pickers by pairwise cut-and-choose. Proposition~\ref{prop:two} certifies each local repair, but we lack a
potential function guaranteeing termination when valuations differ. Proving such a potential, or finding a counterexample
with $n\ge3$ and $m\ge7$ beyond our search range, would resolve the problem.

\section{Completion Estimate}
52\%

This is a post hoc, heuristic estimate for the full catalogue target.
A threshold-graph characterization distinguishes the feasible pair-demand class from additive top-two valuations and supports two-agent, identical-valuation and few-goods existence cases. The general heterogeneous allocation argument still lacks a terminating repair invariant, and computational searches do not prove universal PMMS existence.

\begin{thebibliography}{9}
\bibitem{bmp} J. Byrka, F. Malinka, T. Ponitka, \emph{Probing EFX via PMMS: (Non-)Existence Results in Discrete Fair
Division}, arXiv:2507.14957v2 (2025), Definition 2.6, Sections 6--7, Open Problem 5. \url{https://arxiv.org/html/2507.14957v2}.
\bibitem{ch} V. Chv\'atal, P. L. Hammer, Aggregation of inequalities in integer programming, \emph{Annals of Discrete
Mathematics} 1 (1977), 145--162.
\end{thebibliography}
\end{document}
